\chapter{ผลการทดลองเบื้องต้นหรือระบบต้นแบบ}
\label{chapter:result}

บทนี้นำเสนอผลการพัฒนาระบบต้นแบบ UniAssist AI ผลการสร้างฐานข้อมูล ผลการประเมินและคัดเลือกโมเดลภาษาสำหรับการสกัดข้อมูล ตลอดจนผลการทดสอบการทำงานของแกน Text-to-SQL และฟังก์ชันเสริมต่าง ๆ

\section{ระบบต้นแบบที่พัฒนาขึ้น}

ระบบต้นแบบถูกพัฒนาเป็นเว็บแอปพลิเคชันที่ประกอบด้วยหน้าใช้งานหลัก 3 ส่วน ได้แก่ หน้าแชทถาม-ตอบข้อมูลหลักสูตร หน้าคำนวณเกรด และหน้าข้อมูลทุนการศึกษา พร้อมด้วยแดชบอร์ดสำหรับอาจารย์ที่ปรึกษา

\subsection{หน้าแชทถาม-ตอบข้อมูลหลักสูตร}

หน้าแชทเป็นส่วนติดต่อหลักที่ผู้ใช้พิมพ์คำถามเป็นภาษาไทย ระบบจะแสดงคำตอบที่เรียบเรียงแล้ว พร้อมทั้งแสดงคำสั่ง SQL ที่โมเดลเขียนและตารางผลลัพธ์ดิบ เพื่อให้ตรวจสอบย้อนกลับได้ ดังรูปที่ \ref{fig:chat}

\begin{figure}[ht]
    \centering
    \includegraphics[width=0.95\linewidth]{chat.png}
    \caption{หน้าแชทถาม-ตอบข้อมูลหลักสูตรของระบบ UniAssist AI}
    \label{fig:chat}
\end{figure}

\subsection{หน้าคำนวณเกรดและประเมินความเสี่ยง}

หน้าคำนวณเกรดให้ผู้ใช้วางข้อความผลการเรียนหรือแก้ไขรายวิชาในตาราง ระบบจะคำนวณ GPA รายภาคการศึกษาและสะสม พร้อมจัดสถานะทางการเรียน (เช่น เสี่ยงพ้นสภาพ ภาคทัณฑ์ หรือเข้าเกณฑ์เกียรตินิยม) ดังรูปที่ \ref{fig:gpa}

\begin{figure}[ht]
    \centering
    \includegraphics[width=0.95\linewidth]{gpa.png}
    \caption{หน้าคำนวณเกรดและประเมินสถานะทางการเรียน}
    \label{fig:gpa}
\end{figure}

\subsection{หน้าข้อมูลทุนการศึกษา}

หน้าข้อมูลทุนการศึกษาแสดงรายการทุนพร้อมเงื่อนไขและคุณสมบัติ ดังรูปที่ \ref{fig:scholarship}

\begin{figure}[ht]
    \centering
    \includegraphics[width=0.95\linewidth]{scholarship.png}
    \caption{หน้าข้อมูลทุนการศึกษา}
    \label{fig:scholarship}
\end{figure}

\section{ผลการสร้างฐานข้อมูล}

จากกระบวนการสกัดและนำเข้าข้อมูล ระบบได้ฐานข้อมูลเชิงสัมพันธ์ที่ครอบคลุมทั้ง 4 หลักสูตรของคณะเทคโนโลยีสารสนเทศ ประกอบด้วยรายวิชา 300 รายวิชา แผนการเรียนรายเทอม 583 รายการ ระเบียบข้อบังคับ 108 รายการ และข้อมูลทุนการศึกษา 11 ทุน ตามที่แสดงในตารางที่ \ref{tab:schema} ของบทที่ \ref{chapter:experiment} ฐานข้อมูลนี้ถูกใช้เป็นแหล่งข้อมูลจริงสำหรับการค้นคืนของแกน Text-to-SQL และการคำนวณต่าง ๆ

\section{ผลการประเมินและคัดเลือกโมเดลภาษาสำหรับการสกัดข้อมูล}

ผลการประเมินคุณภาพการสกัดข้อมูลของโมเดลทั้ง 4 ตัว เทียบกับเอกสาร มคอ.2 ฉบับจริง สรุปได้ดังตารางที่ \ref{tab:model-summary}

\begin{table}[ht]
    \centering
    \caption{สรุปผลการเปรียบเทียบความครบถ้วนและความถูกต้องของการสกัดข้อมูลแต่ละโมเดล}
    \label{tab:model-summary}
    \begin{tabularx}{\linewidth}{l c c c X}
        \toprule
        \textbf{โมเดล} & \textbf{Coverage} & \textbf{Accuracy} & \textbf{แอตทริบิวต์} & \textbf{ข้อสังเกต} \\
        \midrule
        Qwen & 91.7\% & 95.2\% & 7/7 & สมดุลดีที่สุด ครบและแม่น \\
        Gemini & 91.7\% & 89.7\% & 7/7 & ครบดี แต่แม่นรองจาก Qwen \\
        DeepSeek & 93.8\% & 84.9\% & 7/7 & ดึงครบที่สุด แต่แม่นน้อยที่สุด \\
        Typhoon & 55.6\% & 97.3\% & 4/7 & แม่นเมื่อดึงเจอ แต่ดึงพลาดเกือบครึ่ง \\
        \bottomrule
    \end{tabularx}
\end{table}

จากตารางที่ \ref{tab:model-summary} จะเห็นว่าแม้โมเดล Typhoon จะมีความถูกต้อง (Accuracy) สูงที่สุด แต่มีความครบถ้วน (Coverage) เพียง 55.6\% และดึงแอตทริบิวต์ได้เพียง 4 จาก 7 รายการ จึงไม่เหมาะกับการสร้างฐานข้อมูลที่ต้องการความครบถ้วน ส่วนโมเดล DeepSeek แม้ดึงรายวิชาได้ครบที่สุด แต่มีความถูกต้องของค่าต่ำที่สุด ขณะที่โมเดล Qwen ให้ผลที่สมดุลที่สุด คือมีความครบถ้วน 91.7\% ควบคู่กับความถูกต้องสูงถึง 95.2\% และดึงครบทั้ง 7 แอตทริบิวต์ ด้วยเหตุนี้จึงเลือกใช้ \textbf{โมเดล Qwen เป็นโมเดลหลัก} ทั้งในกระบวนการสกัดข้อมูลและการใช้งานจริงในระบบ

ตารางที่ \ref{tab:coverage-curriculum} แสดงความครบถ้วนของการสกัดรายวิชาแยกตามหลักสูตร โดยตัวเลขในวงเล็บคือจำนวนรายวิชาทั้งหมด (Union) ของแต่ละหลักสูตร

\begin{table}[ht]
    \centering
    \caption{ความครบถ้วน (Coverage) ของการสกัดรายวิชาแยกตามหลักสูตร}
    \label{tab:coverage-curriculum}
    \begin{tabularx}{\linewidth}{X c c c c}
        \toprule
        \textbf{หลักสูตร (จำนวนวิชา)} & \textbf{Qwen} & \textbf{Gemini} & \textbf{DeepSeek} & \textbf{Typhoon} \\
        \midrule
        IT (110) & 93.6\% & 90.9\% & 94.5\% & 50.0\% \\
        DSBA (90) & 90.0\% & 91.1\% & 93.3\% & 48.9\% \\
        BIT (69) & 85.5\% & 89.9\% & 89.9\% & 62.3\% \\
        AIT (55) & 98.2\% & 96.4\% & 98.2\% & 69.1\% \\
        \bottomrule
    \end{tabularx}
\end{table}

นอกจากนี้ เมื่อพิจารณาความสอดคล้องของค่าในแต่ละแอตทริบิวต์ (เฉลี่ยรวมทุกหลักสูตร) ดังตารางที่ \ref{tab:attr-agreement} พบว่าแอตทริบิวต์ที่โมเดลสกัดได้แม่นยำที่สุดคือชื่อวิชาภาษาไทย (95.4\%) รองลงมาคือประเภทวิชา (86.3\%) ขณะที่แอตทริบิวต์ที่ท้าทายที่สุดคือปี/ภาคการศึกษา (61.5\%) และหมวดวิชา (63.3\%) เนื่องจากข้อมูลส่วนนี้ในเอกสาร มคอ.2 มักอยู่ในรูปแบบที่ตีความได้หลายแบบ เช่น รายวิชาเลือกที่ลงเรียนได้หลายเทอม

\begin{table}[ht]
    \centering
    \caption{ความสอดคล้องของค่ารายแอตทริบิวต์ (เฉลี่ยรวมทุกหลักสูตร)}
    \label{tab:attr-agreement}
    \begin{tabularx}{\linewidth}{X c}
        \toprule
        \textbf{แอตทริบิวต์} & \textbf{ความสอดคล้องเฉลี่ย} \\
        \midrule
        ชื่อวิชา (ภาษาไทย) & 95.4\% \\
        ประเภทวิชา (บังคับ/เลือก) & 86.3\% \\
        ชื่อวิชา (ภาษาอังกฤษ) & 75.6\% \\
        หน่วยกิต & 74.8\% \\
        หมวดวิชา & 63.3\% \\
        ปี/ภาคการศึกษา & 61.5\% \\
        \bottomrule
    \end{tabularx}
\end{table}

\section{ผลการทดสอบแกน Text-to-SQL}

จากการทดสอบการทำงานของแกน Text-to-SQL พบว่าระบบสามารถแปลงคำถามภาษาไทยเป็นคำสั่ง SQL และตอบคำถามข้อมูลหลักสูตรได้อย่างถูกต้อง เช่น คำถามเกี่ยวกับจำนวนหน่วยกิตที่ต้องเรียนจบของแต่ละหลักสูตร รายวิชาในแผนการเรียนแต่ละเทอม รายวิชาบังคับก่อน และเงื่อนไขระเบียบต่าง ๆ โดยระบบแสดงคำสั่ง SQL ที่ใช้ควบคู่กับคำตอบเสมอ ทำให้สามารถตรวจสอบความถูกต้องได้

กลไกตรวจสอบความปลอดภัยของคำสั่งทำงานได้ตามที่ออกแบบไว้ กล่าวคือ คำสั่งที่ไม่ใช่ \texttt{SELECT}/\texttt{WITH} หรือมีคำสั่งเขียน/แก้ไขข้อมูล จะถูกปฏิเสธก่อนรัน และคำสั่งที่ไม่มี \texttt{LIMIT} จะถูกเติมโดยอัตโนมัติ นอกจากนี้ กลไกสำรอง (Fallback) ช่วยให้ระบบยังตอบผู้ใช้ได้แม้ในกรณีที่ไม่สามารถสร้างคำสั่ง SQL ที่เหมาะสม โดยสลับไปโหมดสนทนาที่ยึดข้อเท็จจริงจากฐานข้อมูล

อย่างไรก็ตาม ในระหว่างการทดสอบพบว่าบริการโมเดลภาษาที่เรียกใช้ผ่านผู้ให้บริการภายนอกมีความไม่คงเส้นคงวา (Nondeterminism) เป็นบางครั้ง กล่าวคือ คำถามเดียวกันอาจได้รับคำตอบ \texttt{NO\_ANSWER} ในบางช่วงเวลาแม้จะเป็นคำถามที่ตอบได้จริง ซึ่งเกิดจากการที่ผู้ให้บริการสลับใช้เซิร์ฟเวอร์ที่มีการบีบอัดโมเดล (Quantization) แตกต่างกัน ปัญหานี้ได้รับการบรรเทาด้วยกลไกการลองสร้างคำสั่งใหม่ (Retry) เมื่อได้ผลลัพธ์ \texttt{NO\_ANSWER} และการเลือกใช้โมเดล Qwen ขนาดใหญ่ที่ให้ผลนิ่งกว่า

\section{ผลการทดสอบระบบคำนวณเกรดและแดชบอร์ด}

ระบบคำนวณเกรดสามารถแยกวิเคราะห์ข้อความผลการเรียนที่มีรูปแบบหลากหลาย คำนวณ GPA รายภาคการศึกษาและสะสมได้ถูกต้อง โดยตัดรายวิชาที่ไม่นับหน่วยกิต GPA ออกอย่างเหมาะสม และจัดสถานะทางการเรียนตามเกณฑ์ที่อ้างอิงจากฐานข้อมูลจริง ส่วนแดชบอร์ดสำหรับอาจารย์ที่ปรึกษาสามารถแสดงภาพรวมผลการเรียนเฉลี่ยรายภาคของนักศึกษาในความดูแล และจำแนกกลุ่มความเสี่ยงได้ตามที่ออกแบบไว้ ช่วยให้อาจารย์สามารถระบุนักศึกษาที่ควรได้รับการดูแลเป็นพิเศษได้อย่างรวดเร็ว
